\section*{3.1 Resistant gadgets (author source pp.14--15)}
\begin{originalresult}{Definition 3.1}
A gadget $g:\bits^m\to\bits$ is resistant if for every $i\in[m]$ and $b\in\bits$, the function obtained by fixing the $i$th input to $b$ is not constant. Equivalently, the minimum certificate complexity at inputs is larger than 1. E.g., the XOR function is resistant while the AND function is not.
\end{originalresult}
\begin{originalresult}{Theorem 3.2}
For any function $f:\bits^n\to\bits$ and a resistant gadget $g:\bits^m\to\bits$,
\[\log\DTSize(f\comp g)=\Omega(\DT(f)).\]
\end{originalresult}
In the following proofs, we use $x_1,\ldots,x_n$ to refer to the inputs of function $f$. The inputs to the lifted function $f\comp g$ are denoted by $z_{i,j}$, where $i\in[n]$ and $j\in[m]$.
\begin{proof}[Proof using random projections]
Let $T$ be a decision tree for $f\comp g$ of size $S$. We apply the following random projection. For each $i\in[n]$, choose an index $j\in[m]$ at random, and set $z_{i,j}$ to a random value $b\in\bits$. Set the remaining variables $z_{i,k}$, $k\ne j$, to $x_i,\bar x_i,0,1$ in such a way that $g(z_{i,1},\ldots,z_{i,m})=x_i$. This is possible since $g$ is resistant, as we show below. The result is a decision tree $T'$ for $f$.
To see that we can always construct such a projection, let $h(z)$ be a non-constant function, and choose two inputs $w^0,w^1$ such that $h(w^b)=b$. We define the projection $\rho$ as follows: if $w_i^0=w_i^1=a$ then set $z_i=a$; otherwise, if $w_i^0=0$ then set $z_i=x$, and if $w_i^1=0$ then set $z_i=\bar x$. By construction, $\rho|_{x=b}=w^b$ and so $h|_\rho(b)=h(w^b)=b$.
For example, suppose that $h(z_1,z_2,z_3)$ is the majority function on three inputs and that we set the first bit $z_1$ to 0. Then $h(0,1,x_i)=x_i$ and $h(0,x_i,x_i)=x_i$ are two possible ways to choose a pair of inputs $(w^0,w^1)$ corresponding to $(010,011)$ and $(000,011)$, respectively.
We would like to say that there is a choice of random projection for which $T'$ is shallow. To do this, let us calculate the probability that some vertex $v$ at depth $D$ in $T$ survives in $T'$. The path leading to $v$ queries variables from at least $D/m$ different blocks. For each such block, choose an arbitrary variable that is queried along the path and consider the value the path assigns to it. If the random projection sets this variable to the opposite value then the vertex $v$ does not survive in $T'$. It happens with probability at least $1/(2m)$ for each variable, and so the vertex $v$ survives with probability $p\le(1-1/(2m))^{D/m}=1/2^{O(D)}$.
If $pS<1$, then the expected number of vertices at depth $D$ surviving in $T'$ is less than 1, and so there is a random projection for which $T'$ has a depth less than $D$. This happens for some $D=O(\log S)$.
\end{proof}
